%% preamble.tex — Category Theory: A Structural Approach
%% Shared preamble for all chapters and appendices

%% ============================================================
%% PACKAGES
%% ============================================================

%% Mathematics
\usepackage{amsmath}
\usepackage{amssymb}
\usepackage{amsthm}
\usepackage{mathtools}      % extends amsmath

%% Graphics and diagrams
\usepackage{tikz}
\usepackage{tikz-cd}        % commutative diagrams
\usetikzlibrary{arrows.meta, decorations.pathmorphing, positioning, calc, matrix}

%% Layout and typography
\usepackage[margin=1.25in]{geometry}
\usepackage{microtype}      % better justification
\usepackage{graphicx}
\usepackage{booktabs}       % better tables
\usepackage{enumitem}       % fine-grained list control

%% Colors and boxes
\usepackage{xcolor}
\usepackage{tcolorbox}
\tcbuselibrary{skins, breakable}

%% References and hyperlinks
\usepackage[colorlinks=true,
            linkcolor=blue!60!black,
            citecolor=green!50!black,
            urlcolor=blue]{hyperref}
\usepackage{cleveref}

%% Bibliography
\usepackage[backend=biber,
            style=alphabetic,
            maxnames=4,
            doi=false,
            isbn=false]{biblatex}
\addbibresource{bibliography.bib}

%% ============================================================
%% TIKZ-CD SETTINGS
%% ============================================================

\tikzcdset{
  arrow style=math font,
  diagrams={>=Straight Barb},
}
\tikzset{
  every node/.style={font=\small},
}

%% ============================================================
%% THEOREM ENVIRONMENTS
%% ============================================================

\theoremstyle{plain}
\newtheorem{theorem}{Theorem}[chapter]
\newtheorem{proposition}[theorem]{Proposition}
\newtheorem{lemma}[theorem]{Lemma}
\newtheorem{corollary}[theorem]{Corollary}

\theoremstyle{definition}
\newtheorem{definition}[theorem]{Definition}
\newtheorem{example}[theorem]{Example}

\theoremstyle{remark}
\newtheorem{remark}[theorem]{Remark}
\newtheorem{warning}[theorem]{Warning}
\newtheorem{exercise}[theorem]{Exercise}

%% ============================================================
%% TCOLORBOX STYLES
%% ============================================================

%% Concept box (light blue, for key concepts / overviews)
\tcbset{
  conceptbox/.style={
    enhanced,
    colback=blue!7!white,
    colframe=blue!40!black,
    fonttitle=\bfseries,
    breakable,
    before upper={\parindent1.5em},
  }
}

%% ============================================================
%% CUSTOM MATH MACROS
%% ============================================================

%% -- Category names ------------------------------------------
\newcommand{\cat}[1]{\mathbf{#1}}          % generic category

\newcommand{\Set}{\mathbf{Set}}
\newcommand{\Grp}{\mathbf{Grp}}
\newcommand{\Top}{\mathbf{Top}}
\newcommand{\Ab}{\mathbf{Ab}}
\newcommand{\Ring}{\mathbf{Ring}}
\newcommand{\Pos}{\mathbf{Pos}}
\newcommand{\Cat}{\mathbf{Cat}}
\newcommand{\Man}{\mathbf{Man}}
\newcommand{\Vect}{\mathbf{Vect}}

%% -- Morphism operations -------------------------------------
\newcommand{\op}{^{\mathrm{op}}}           % opposite category
\newcommand{\id}{\mathrm{id}}              % identity morphism
\newcommand{\Id}{\mathrm{Id}}              % identity functor
\newcommand{\iso}{\xrightarrow{\sim}}      % isomorphism arrow
\newcommand{\blank}{{-}}                   % placeholder

%% -- Hom and related ----------------------------------------
\DeclareMathOperator{\Hom}{Hom}
\DeclareMathOperator{\End}{End}
\DeclareMathOperator{\Aut}{Aut}
\DeclareMathOperator{\Nat}{Nat}
\DeclareMathOperator{\Fun}{Fun}
\DeclareMathOperator{\ob}{ob}
\DeclareMathOperator{\mor}{mor}
\DeclareMathOperator{\dom}{dom}
\DeclareMathOperator{\cod}{cod}

%% -- Limits and colimits ------------------------------------
\newcommand{\colim}{\operatorname{colim}}
\newcommand{\holim}{\operatorname{holim}}
\newcommand{\hocolim}{\operatorname{hocolim}}

%% -- Kan extensions -----------------------------------------
\newcommand{\Lan}{\operatorname{Lan}}
\newcommand{\Ran}{\operatorname{Ran}}

%% -- Adjunction symbol --------------------------------------
\newcommand{\adj}{\mathrel{\dashv}}        % F ⊣ G

%% -- Yoneda embedding ---------------------------------------
\newcommand{\yo}{\mathbf{y}}

%% -- Tensor product -----------------------------------------
\newcommand{\tensor}{\otimes}

%% -- Number systems ------------------------------------------
\newcommand{\NN}{\mathbb{N}}
\newcommand{\ZZ}{\mathbb{Z}}
\newcommand{\QQ}{\mathbb{Q}}
\newcommand{\RR}{\mathbb{R}}
\newcommand{\CC}{\mathbb{C}}

%% -- Misc ---------------------------------------------------
\newcommand{\defeq}{\coloneqq}            % :=
\newcommand{\eqdef}{\eqqcolon}            % =:

%% ============================================================
%% CLEVEREF NAMES
%% ============================================================

\crefname{theorem}{Theorem}{Theorems}
\crefname{proposition}{Proposition}{Propositions}
\crefname{lemma}{Lemma}{Lemmas}
\crefname{corollary}{Corollary}{Corollaries}
\crefname{definition}{Definition}{Definitions}
\crefname{example}{Example}{Examples}
\crefname{remark}{Remark}{Remarks}
\crefname{exercise}{Exercise}{Exercises}
\crefname{chapter}{Chapter}{Chapters}
\crefname{appendix}{Appendix}{Appendices}
\crefname{section}{Section}{Sections}

%% ============================================================
%% MISC SETTINGS
%% ============================================================

%% Wider display math
\setlength{\jot}{6pt}

%% Penalty settings to reduce widow/orphan lines
\widowpenalty=10000
\clubpenalty=10000
